\section{Integral attachment and cochain descent}\label{sec:integral}

The useful reduction in the middle degree is elementary once the
correct integral restrictions have been calculated. We state it
first, so that the later Grassmannian computation has a definite
target: the homomorphism \(\kappa\), not just its finite codomain.

\subsection{A calculated boundary homomorphism}
Let \(M\) be a closed oriented smooth \(2n\)-manifold and \(E\subset M\)
a compact tame subset. Suppose it has a closed oriented regular
neighborhood \(W\) retracting onto \(E\), with smooth boundary \(L\).
Put \(V=M\setminus E\). Collars identify the overlap of \(V\) with an
open neighborhood of \(E\), up to homotopy, with \(L\).
Assume
\begin{equation}\label{eq:attachment-assumptions}
 \begin{gathered}
 H_n(E,\Z)=\Lambda\ \text{is free},\qquad H^n(E,\Z)=\Lambda^*,\\
 H^{n-1}(E,\Z)=H_{n-1}(E,\Z)=0,\qquad
 B:\Lambda\longrightarrow\Lambda^*\ \text{is injective}.
 \end{gathered}
\end{equation}
Here the dual identification is evaluation, and \(B\) is the middle
intersection map in \(W\). Let \(C=\coker B\).

\begin{theorem}\label{thm:attachment}
Suppose homogeneous integral classes \(g_j\in H^*(M,\Z)\) have
degree-\(n\) monomials generating \(H^n(M,\Z)\).
Supply the ring \(A=H^*(V,\Z)\) with the distinguished classes
\(a_j=g_j|_V\), and calculate \(r_j=g_j|_E\).
Then \(H^{n-1}(L,\Z)=0\), \(H^n(L,\Z)=C\), and
\begin{equation}\label{eq:kappa}
 \kappa\bigl(P(a_j)\bigr)=[P(r_j)]\in C
\end{equation}
is a well-defined homomorphism on all of \(A^n\).
It is the actual boundary restriction and gives
\begin{equation}\label{eq:fibre-product}
 H^n(M,\Z)\ \simeq\
 \{(u,z)\in A^n\oplus\Lambda^*:\ \kappa(u)=[z]\text{ in }C\}.
\end{equation}
The isomorphism is restriction to \(V\) and \(W\).
\end{theorem}

\begin{proof}
Poincaré--Lefschetz duality and the sequence of \((W,L)\) give
\[
 0\longrightarrow H^{n-1}(L,\Z)\longrightarrow\Lambda
 \xrightarrow{B}\Lambda^*\longrightarrow H^n(L,\Z)
 \longrightarrow0.
\]
This proves the boundary assertions. Excision also gives
\(H^{n+1}(M,V;\Z)=H_{n-1}(E,\Z)=0\).
Therefore restriction \(H^n(M,\Z)\to A^n\) is onto, and the
monomials in the \(a_j\) generate \(A^n\).

It remains to check every relation, which is the point of the
formula. If \(P(a_j)=0\), the class \(P(g_j)\) comes from
\(H^n(M,V;\Z)=\Lambda\). Its restriction to \(W\) is the image of
that same element under \(B\). Hence \(P(r_j)\in B\Lambda\),
so \eqref{eq:kappa} is independent of its polynomial representative.
The two restrictions of each \(g_j\) agree on \(L\); consequently
\(\kappa\) is the boundary restriction. The Mayer--Vietoris sequence,
with \(H^{n-1}(L,\Z)=0\), now identifies \(H^n(M,\Z)\) with the
kernel of \((u,z)\mapsto\kappa(u)-[z]\).
\end{proof}

This is an integral fibre product. It retains the finite boundary
group, saturation and extension information that would disappear
after tensoring with \(\QQ\). It requires the marked complement ring;
it does not calculate that ring from the singularity. A basis change
in \(\Lambda\) gives the naturally isomorphic fibre product.
Changing the generating set changes the formula, but not the
homomorphism, because it is the same boundary restriction.
Isotopic regular neighborhoods give the same result under collar
identifications.

\subsection{Integral generators from universal indices}
For primitive-vector K3 moduli, the required global generators come
from \eqref{eq:mukai-convention}. Markman's theorem
\cite[Theorem~1]{markman} states that Chern classes of the topological
\(K\)-theory Künneth factors of a universal class generate integral
cohomology. The integral Euler pairing on \(K^0_{\rm top}(X)\) is
perfect. Taking an Euler-dual integral basis therefore expresses
those factors as the indices \(I_z\). Thus the classes \(c_k(I_z)\)
generate \(H^*(M,\Z)\). The same theorem gives odd vanishing and
torsion-freeness in this K3 setting.

Both qualifications matter: all topological \(K\)-classes are
allowed, and the generation is integral. Algebraic indices alone,
or rational Chern characters alone, would not justify
Theorem~\ref{thm:attachment}. In a concrete exceptional family,
a base-line ambiguity can often be removed by the determinant of a
rank-one global index. Section~\ref{sec:schubert} performs exactly
this normalization before calculating \(\kappa\).

\subsection{All degrees and all overlaps}\label{subsec:cochains}
The general construction uses the geometric data of
Section~\ref{subsec:retained}. Whole-fibre coverage, proved in
Proposition~\ref{prop:coverage}, supplies a finite open cover
\(\mathcal U=(U_i)\) of \(M_\theta\) by those domains, together with
domains in the regular locus. Write \(U_I=\bigcap_{i\in I}U_i\).
Every intersection is expressed in any one of its coordinates by
the actual classifying transitions.

\begin{proposition}\label{prop:descent}
Let
\[
 K^{p,q}(\mathcal U)=
 \prod_{i_0<\cdots<i_p} C^q(U_{i_0\cdots i_p};\Z).
\]
With alternating restriction differential \(\delta\) and total
differential \(D=\delta+(-1)^p d\), restriction induces a
quasi-isomorphism
\begin{equation}\label{eq:cech}
 C^*(M_\theta;\Z)\longrightarrow\operatorname{Tot}K(\mathcal U).
\end{equation}
Changing the cover, slices or family trivializations gives the same
object in the derived category of integral complexes.
\end{proposition}

\begin{proof}
The coarse coordinate maps are unique on their overlaps and hence
satisfy the cocycle identity on triples, and all its higher
consequences. Thus \(D^2=0\).

Use first the singular chains subordinate to \(\mathcal U\).
Barycentric subdivision and its usual chain homotopy show that their
inclusion in all singular chains is a chain-homotopy equivalence.
For each subordinate simplex, the set of indices of cover members
containing its image is a nonempty simplex. Its augmented Čech
complex is contractible. Taking the corresponding products proves
exactness of the augmented cochain resolution in the Čech direction.
The Čech degree is bounded because the cover is finite; its
filtration gives \eqref{eq:cech}.

A refinement commutes with augmentation, so its total-complex map
is a quasi-isomorphism. Two covers have a common intersection
refinement. Different refinement assignments induce the same map
in the derived category, since both commute with the invertible
augmentation from \(C^*(M_\theta;\Z)\). Finally a change of slices
or trivializations conjugates the overlap diagram by the unique
coarse coordinate maps. Base-line twists, and different presentations
of a universal gerbe, do not change those coarse maps.
\end{proof}

The descent assertion is classical; compare the open-cover
construction in \cite[Section~4G]{hatcher}. Its application here uses the
geometric identification of the domains and their overlaps, rather
than replacing them by lists of fibre cohomology groups.
The complex retains all multiple intersections and monodromy.
It computes additive integral cohomology; no cup-product or
homotopy-type assertion is being inferred from this additive
statement. A finite cover gives a finite diagram of complexes,
not automatically a finite matrix calculation. Compatible cellular
models can make it finite when their maps are also supplied.

In particular, the formal descent argument does not solve the
existence or compatibility problem for arbitrary local resolutions.
The local-to-global theory of Kaplan--Schedler
\cite{kaplanschedler} treats such compatibility and monodromy
questions in a broader setting. Here the compact projective
chamber spaces are already part of (G1). Our task is to identify
whole local fibres within them and to calculate the integral
attachments on their universal-index classes.
